\section{THEORETICAL FRAMEWORK OF TMDs}

\subsection{Definition of TMDs} 
The theoretical foundation for accessing Transverse Momentum Dependent Parton Distribution Functions (TMDs) relies on analyzing a quark-quark correlation function. We define the general unsubtracted correlator as:
\begin{align}
	\Phi^{[\Gamma]}(\kei,P,S;\ldots)\ \equiv\ \ 
	\int \frac{d^4 \bvec}{(2\pi)^4}\, e^{i \kei\cdot \bvec}\,
		\frac{\overbrace{\rule{0em}{1.2em}
		\frac{1}{2}\, \bra{P,S}\ \bar{q}(0)\, \Gamma\ \WlineC{\mathcal{C}_\bvec}\ q(\bvec)\ \ket{P,S}
		}^{\displaystyle \equiv \widetilde \Phi_\unsub^{[\Gamma]}(\bvec,P,S;\ldots)}}
		{ \widetilde{\softf}(\bvec^2;\ldots) } \ .
		\label{eq-corr}	
\end{align}
The numerator, $\widetilde \Phi_\unsub^{[\Gamma]}(\bvec,P,S;\ldots)$, captures the bare non-local matrix element where the quark fields are connected by a gauge link $\WlineC{\mathcal{C}_\bvec}$. The denominator introduces the soft factor, $\widetilde{\softf}(\bvec^2;\ldots)$, which is essential for regulating rapidity divergences. The ellipses (``$\ldots$'') denote additional dependence on parameters governing the Wilson line geometry and the soft factor, as well as an implicit dependence on the ultraviolet renormalization scale $\mu$.

Depending on the specific factorization scheme employed, the soft factor $\widetilde{\softf}$ is constructed from vacuum expectation values of Wilson lines. Following the framework established by Collins and Aybat \cite{Aybat:2011zv,CollinsBook2011,Aybat:2011ge}, this takes the specific form:
\begin{equation}
\widetilde{\softf}(\bvec^2;\ldots) = 
  \sqrt{ \frac
    { \tilde S_{(0)}( \vprp{\bvec}, + \infty, -\infty )\ \tilde S_{(0)}( \vprp{\bvec}, y_s, -\infty ) }
    { \tilde S_{(0)}( \vprp{\bvec}, + \infty, y_s ) }
  } \ ,
\end{equation}
where $\tilde S_{(0)}( \vprp{\bvec}, \ldots )$ represents these vacuum expectation values. In this scheme, the Wilson lines are initially space-like, with some remaining tilted off the light cone (characterized by the rapidity $y_s$). The lines within $\WlineC{\mathcal{C}_\bvec}$ in the numerator of Eq.~\eqref{eq-corr} are formally taken to the light-cone limit ($\pm\infty$). While space-like Wilson line structures are directly computable on the Euclidean lattice, executing the strict light-cone limit numerically remains a profound challenge. However, for the specific observables constructed in this work, the explicit functional form of $\widetilde{\softf}(\bvec^2;\ldots)$ does not need to be rigorously evaluated, as it strategically cancels out in appropriately constructed ratios.

Projecting the correlator onto the relevant kinematic space requires integrating over the suppressed momentum component $k^-$, which corresponds to evaluating the position-space separation at $b^+ = 0$:
\begin{align}
	\Phi^{[\Gamma]}(x,\vprp{k};P,S;\ldots) & \equiv \int dk^{-} \Phi^{[\Gamma]}(\kei,P,S;\ldots) \nonumber \\ 
	& =\int \frac{d^2 \vprp{\bvec}}{(2\pi)^2} \int \frac{d (\bvec \tcdot P)}{(2\pi) P^+}\ e^{i x (\bvec \tcdot P) - i\vprp{\bvec} \tcdot \vprp{\kei}}\,
		\left. \frac{\frac{1}{2}\, \bra{P,S}\ \bar{q}(0)\, \Gamma\ \WlineC{\mathcal{C}_\bvec}\ q(\bvec)\ \ket{P,S}}
		{ \widetilde{\mathcal{S}}(-\vprp{\bvec}^2;\ldots) } \right|_{b^+ = 0} \ .
\end{align}
In this regime, $x$ and $(\bvec \tcdot P)$ form a Fourier conjugate pair, as do $\vprp{k}$ and $\vprp{b}$. By selecting the Dirac matrix $\Gamma = \gamma^+$ to project onto the leading-twist unpolarized quark distribution, the correlator yields \cite{Ralston:1979ys,Tangerman:1994eh,Mulders:1995dh,Goeke:2005hb}:
\begin{align}
	\Phi^{[\gamma^+]}(x,\vprp{k};P,S,\ldots) & = f_1 - \toddmark{\frac{\myeps_{ij}\, \vect{k}_{i}\, \vect{S}_{j}}{m_N}\ f_{1T}^\prp} \label{eq-phigammaplus}\ . 
\end{align}
The term enclosed in $[\ ]_\text{odd}$ isolates the naively time-reversal odd (T-odd) Sivers function $f_{1T}^\perp$ \cite{Sivers:1989cc}, representing the distribution of unpolarized quarks inside a transversely polarized nucleon. The emergence and survival of this T-odd structure is inextricably linked to the extended gauge link geometry and symmetry transformations discussed in the subsequent sections.

\subsection{General strategy}
Transitioning this continuous framework onto a discrete Euclidean spacetime lattice introduces profound theoretical constraints. Evaluating standard collinear PDFs utilizes a straightforward, light-like Wilson line:
\begin{equation*}
  f_1(x) \equiv \frac{1}{2 (2\pi)} \int db^- e^{i x P^+ b^-} \, \bra{P,S}\ \bar{q}(0)\, \gamma^+\ \Wline{0,n b^-}\ q(n b^-)\ \ket{P,S} \ .
\end{equation*}
Because this interval is strictly Minkowskian, it cannot be rotated into Euclidean space, necessitating an Operator Product Expansion (OPE) approach to reduce the non-local separation into local matrix elements \cite{Collins:1981uw}. 

Extracting the $x$-dependence of the Sivers TMD, however, deviates from this collinear constraint and presents unique lattice opportunities:
\begin{enumerate}
  \item \textbf{Space-like Transverse Separation:} The operator separation acquires a transverse component, $\bvec = n \bvec^- + \bvec_\perp$. This ensures the interval is generally space-like, permitting a direct representation of the non-local operator in Euclidean spacetime.
  \item \textbf{Extended Gauge Link Geometries:} The Wilson line $\WlineC{\mathcal{C}_\bvec}$ is process-dependent and typically extends to infinity to model initial or final state interactions. This complex topology fundamentally precludes a simple expansion into local operators. 
  \item \textbf{Soft Factor Regularization:} Standard $\overline{\mathrm{MS}}$ dimensional regularization fails to capture the rapidity divergences associated with these extended links, forcing the inclusion of the soft factor $\tilde \softf$ and mandating evolution across additional regularization scales.
\end{enumerate}

\subsection{TMDs in Fourier space}
Because our lattice methodology natively extracts the $b$-dependent matrix elements $\widetilde \Phi_\unsub^{[\Gamma]}(\bvec,P,S;\ldots)$, evaluating the Sivers observable directly in Fourier space ($\vprp{b}$-space) is phenomenologically preferred. For a generic TMD $f$, we define the Fourier transform and its derivatives:
\begin{align}
	\tilde f(x, \vprp{\bvec}^2;\ldots) & \equiv \int  d^2 \vprp{\kei} \, e^{i \vprp{\bvec} \cdot \vprp{\kei} }\; f(x, \vprp{\kei}^2;\ldots)
		= 2\pi \int      d|\vprp{\kei}| |\vprp{\kei}|\ J_0(|\vprp{\bvec}||\vprp{\kei}|)\  f(x, {\vprp{\kei}^2};\ldots )\ , \label{eq-FTMD} \\
	\tilde f^{(n)}(x, \vprp{\bvec}^2\ldots ) & \equiv n!\left( -\frac{2}{\mN^2}\partial_{\vprp{\bvec}^2} \right)^n \ 
	\tilde f(x, \vprp{\bvec}^2;\ldots ) 
		= \frac{ 2\pi \ n!}{(\mN^2)^n} \int  d |\vprp{\kei}| |\vprp{\kei}|\left( \frac{|\vprp{\kei}|}{|\vprp{\bvec}|}\right)^n J_n(|\vprp{\bvec}||\vprp{\kei}|)\  f(x, {\vprp{\kei}^2};\ldots )\ , \label{eq-dFTMD}
\end{align}
where $J_n$ are the Bessel functions of the first kind. As articulated in Ref.~\cite{Boer:2011xd}, these Bessel-weighted quantities provide a robust pathway for structuring evolution equations. 

Taking the strict $|\vprp{\bvec}| \rightarrow 0$ limit theoretically recovers the conventional $\vprp{\kei}$-moments:
\begin{align}
	\tilde f^{(n)}(x, 0;\ldots ) & = \int  d^2 \vprp{\kei} \left( \frac{\vprp{\kei}^2}{2 \mN^2 }\right)^n  f(x, \vprp{\kei}^2;\ldots ) \equiv  f^{(n)}(x)\ .
	\label{eq-btder2}
\end{align}
However, the first transverse moment of the Sivers function, $f_{1T}^{\perp(1)}(x)$, is notoriously divergent without additional regularization, as the perturbatively calculable high-$\vprp{\kei}$ tails do not decay fast enough \cite{Bacchetta:2008xw}. To circumvent this divergence, we avoid the $\vprp{\bvec}=0$ extrapolation entirely. By retaining finite $|\vprp{\bvec}|$, the integrals remain convergent. The corresponding $x$-dependence of the Sivers TMD is then accessed by mapping the conjugate variable $\bvec \tcdot P$ \cite{Hagler:2009mb,Musch:2010ka}.

\subsection{Link geometry}
To accommodate the symmetries required to isolate the Sivers effect, we employ a staple-shaped gauge link:
\begin{align}
	\WlineC{\mathcal{C}_\bvec^{(\eta v)}} = \Wline{0, \eta v, \eta v + b, b}\ .
	\label{eq-staplelink}
\end{align}
The vector $v$ governs the direction of the staple, and $\eta$ controls its extent. By ensuring $v$ is space-like ($\vprp{v}=0$), we maintain the ability to simulate this operator on a Euclidean lattice \cite{Ji:2004wu,Aybat:2011zv,CollinsBook2011}. Focusing on the lowest $x$-moment ($b^- = b^+ = 0$), the link at $\eta v$ acts as a purely transverse connector.

The approach toward the physical light-like limit ($v = \nminus$) is parameterized by the Lorentz-invariant Collins-Soper parameter, $\zeta \equiv 2 v \tcdot P/\sqrt{|v^2|}$. As $v$ becomes light-like, $\zeta \rightarrow \infty$, acting as an artificial scale regulating the rapidity divergences \cite{Collins:1981uk, Idilbi:2004vb, Aybat:2011ge}. Expressing the vectors $P$ and $v$ via rapidities ($y_P$ and $y_v$), we define the dimensionless variable:
\begin{equation}
	\hat \zeta \equiv \zeta / 2 \mN =  \frac{v \cdot P}{\sqrt{|v^2|} \sqrt{P^2}} = \sinh(y_P - y_v)\ ,
	\label{eq-zeta}
\end{equation}
highlighting that $\hat \zeta$ effectively measures the rapidity difference between the nucleon momentum and the Wilson line direction. Crucially for our methodology, this permits boosting the system to a frame where $v$ is purely spatial ($v^0 = 0, y_v=0$).

\subsection{Symmetry Constraints}
The correlation function is severely restricted by discrete symmetries. Consider the unsubtracted correlator \cite{Boer:2011xd}:
\begin{align}
  \Phi^{[\GammaOp]}_{\text{unsub}} (\kei,P,S;v,\mu) & = \int \frac{d^4 \elll}{(2\pi)^4} \ 
  e^{i\kei \cdot \elll} 
  \underbrace{ \frac{1}{2} \bra{\nucl{P,S}}\ \bar \quark(0)\ \overbrace{ \mathcal{U}[0,\infty v]\ \mathcal{U}[\infty v,\elll] }^{\displaystyle \U} \GammaOp\ \quark(\elll)\ \ket{\nucl{P,S}} }_{\displaystyle \widetilde \Phi^{[\GammaOp]}_{\text{unsub}}(\elll,P,S;v,\mu) }\ .
  \label{eq-corrunmod}
\end{align}
Decomposing this into real-valued Lorentz-invariant amplitudes for the unpolarized quark fields ($\Gamma = \gamma^\mu$) yields:
\begin{align}
\frac{1}{2} \Phi^{[\gamma^\mu]}_{\text{unsub}} & = P^\mu\, \Amp_2 +\kei^\mu\, \Amp_3 + \frac{1}{\mN} \epsilon^{\mu \nu \alpha \beta} P_\nu \kei_\alpha S_\beta\, \Amp_{12} + \frac{\mN^2}{(v \tcdot P)} v^\mu\, \Bmp_1 \nonumber\\
& + \frac{\mN}{v \tcdot P} \epsilon^{\mu \nu \alpha \beta} P_\nu v_\alpha S_\beta\, \Bmp_7 + \frac{ \mN}{v \tcdot P} \epsilon^{\mu \nu \alpha \beta} \kei_\nu v_\alpha S_\beta\, \Bmp_8 \nonumber\\ 
& + \frac{1}{\mN (v \tcdot P)} (\kei \tcdot S) \epsilon^{\mu \nu \alpha \beta} P_\nu \kei_\alpha v_\beta\, \Bmp_9  + \frac{\mN}{(v \tcdot P)^2} (v \tcdot S) \epsilon^{\mu \nu \alpha \beta} P_\nu \kei_\alpha v_\beta \Bmp_{10}\ . 
\label{eq-phidecomp}
\end{align}
Enforcing Lorentz ($L$), Parity ($P$), Time-reversal ($T$), and Hermitian conjugation ($\dagger$) transformations imposes the following constraints:
\begin{align}
	&(L): &\Phi^{[\Gamma]}_{\text{unsub}}(k,P,S;v,\mu)  
	&= \Phi^{[\Lambda_{1/2}^{-1}\Gamma\Lambda_{1/2}^{\phantom{-1}}]}_{\text{unsub}}(\Lambda k,\Lambda P, \Lambda S;\Lambda v,\mu) \ ,\label{eq-lortrans} \displaybreak[0] \\
	&(P): &\Phi^{[\Gamma]}_{\text{unsub}}(k,P,S;v,\mu)  
	&= \Phi^{[\gamma^0\Gamma\gamma^0]}_{\text{unsub}}(\overline{k},\overline{P},-\overline{S};\overline{v},\mu) \ ,\label{eq-conpar} \displaybreak[0] \\
	&(T): &\left[ \Phi^{[\Gamma]}_{\text{unsub}}(k,P,S;v,\mu) \right]^* &= \Phi^{[\gamma^1 \gamma^3 \Gamma^* \gamma^3 \gamma^1]}_{\text{unsub}}(\overline{k},\overline{P},\overline{S};-\overline{v},\mu) \ , \label{eq-contime} \displaybreak[0] \\
	&(\dagger): &\left[  \Phi^{[\Gamma]}_{\text{unsub}}(k,P,S;v,\mu) \right]^* &= \Phi^{[\gamma^0 \Gamma^\dagger \gamma^0]}_{\text{unsub}}(k,P,S;v,\mu) \ . \label{eq-conherm}
\end{align}
While hermiticity forces $\Amp_i$ and $\Bmp_i$ to be real, time reversal flips the sign of $v \tcdot P$, directly bridging the amplitudes governing Semi-Inclusive Deep Inelastic Scattering (SIDIS) and the Drell-Yan (DY) process—the physical mechanism driving the Sivers sign change.

Mapping these systematic restrictions to coordinate space under hermitian conjugation specifically, the position-space correlator transforms as:
\begin{align}
	&(\dagger): &\left[ \widetilde{\Phi}^{[\GammaOp]}_{\text{unsub}}(\elll,P,S;v)  \right]^* &= \widetilde{\Phi}^{[\gamma^0 \GammaOp^\dagger \gamma^0]}_{\text{unsub}}(-\elll,P,S;v)  \ .\label{eq-conhermtilde}
\end{align}
Substituting $\kei \rightarrow -i \mN^2 \elll$, the coordinate parameterization of $\widetilde{\Phi}$ is explicitly constructed. The amplitudes $\tAmp_i$ and $\tBmp_i$ in this space are complex, where their real and imaginary parts strictly map to $b$-even and $b$-odd functions, respectively:
\begin{align}
	\left[ \tAmp_i(\elll^2,\elll \tcdot P, v \tcdot \elll / (v \tcdot P), \zeta^{-2},\mu^2) \right]^* = \tAmp_i(\elll^2,-\elll \tcdot P, -v \tcdot \elll / (v \tcdot P), \zeta^{-2},\mu^2)\, .\label{conditi_hermitian}
\end{align}
\begin{align}
    \tilde{A}^{\text{Re}}_i&=\tilde{A}^{\text{(b-even)}}_i, \quad  \tilde{A}^{\text{Im}}_i=\tilde{A}^{\text{(b-odd})}_i \\
    \tilde{B}^{\text{Re}}_i&=\tilde{B}^{\text{(b-even)}}_i, \quad  \tilde{B}^{\text{Im}}_i=\tilde{B}^{\text{(b-odd})}_i
\end{align}

\subsection{Parametrization of the correlator}
\label{sec-parametr}
Extracting the $x$-dependence of the Sivers TMD from Euclidean data requires a rigorous mapping between the discretized observables and the manifestly Lorentz-invariant amplitudes. Building on the infinite-link parameterization of Goeke et al. \cite{Goeke:2005hb}, and identifying $k \rightarrow - i \mN^2 \bvec$ \cite{Boer:2011xd}, the unpolarized $b$-dependent correlator yields:
\begin{align}
	\frac{1}{2}\widetilde \Phi^{[\gamma^\mu]}_{\unsub} & = 
		P^\mu\, \tAmp_2 - i \mN^2 \bvec^\mu\, \tAmp_3
		- i \mN \epsilon^{\mu \nu \alpha \beta} P_\nu \bvec_\alpha S_\beta\, \tAmp_{12} \nonumber \\ &
		+ \frac{\mN^2}{(v \tcdot P)} v^\mu\, \tBmp_1 
		+ \frac{\mN}{v \tcdot P} \epsilon^{\mu \nu \alpha \beta} P_\nu v_\alpha S_\beta\, \tBmp_7  
		- \frac{ i \mN^3}{v \tcdot P} \epsilon^{\mu \nu \alpha \beta} \bvec_\nu v_\alpha S_\beta\, \tBmp_8 \nonumber \\ &
		- \frac{\mN^3}{v \tcdot P} (\bvec \tcdot S) \epsilon^{\mu \nu \alpha \beta} P_\nu \bvec_\alpha v_\beta\, \tBmp_9
		- \frac{i \mN^3}{(v \tcdot P)^2} (v \tcdot S) \epsilon^{\mu \nu \alpha \beta} P_\nu \bvec_\alpha v_\beta \tBmp_{10}
		\label{eq-phitildedecomp-vec} \displaybreak[0] 
		 \ .
\end{align}
While prior lattice campaigns utilized straight gauge links—which isolated only T-even structures \cite{Hagler:2009mb,Musch:2010ka}—the explicit inclusion of the staple direction $v$ enables the analysis of T-odd observables like the Sivers shift. 

To conduct calculations at finite staple extents on the lattice, we redefine the scale-invariant infinite-link amplitudes to explicitly track the extent $\eta$, mapping them via:
\begin{align}
	\tAmp_i\left(\bvec^2,\bvec \tcdot P,\frac{v \tcdot \bvec}{v \tcdot P}, \frac{v^2}{(v \tcdot P)^2}, \eta v \tcdot P\right) & = \tilde a_i(\bvec^2,\bvec \tcdot P, \eta v \tcdot b, (\eta v)^2, \eta v \tcdot P)   \, , \nonumber \\
	 \tBmp_i\left(\bvec^2,\bvec \tcdot P,\frac{v \tcdot \bvec}{v \tcdot P}, \frac{v^2}{(v \tcdot P)^2}, \eta v \tcdot P\right) & = (\eta v \tcdot P)^n\ \tilde b_i(\bvec^2,\bvec \tcdot P, \eta v \tcdot b, (\eta v)^2, \eta v \tcdot P)   \, .
	\label{eq-ab}
\end{align}
In our frame choice ($b^+=0$, $\vprp{v}=\vprp{P}=0$), kinematic constraints enforce a rigid relationship between the invariants \cite{MuschThesis2,Accardi:2009au}:
\begin{align}
	\frac{v \tcdot b}{v \tcdot P} = b \tcdot P \frac{R(\zetahat^2)}{\mN^2}\, , \quad \text{where} \quad R(\zetahat^2) \equiv 1- \sqrt{1 + \hat \zeta^{-2}} = \frac{\mN^2}{v \tcdot P} \frac{v^+}{P^+} \, .
	\label{eq-lindeprel}
\end{align}
Focusing exclusively on the $\gamma^+$ projection necessary to isolate the unpolarized quark distribution, we combine the amplitudes tailored for Sivers extraction:
\begin{align}
	\frac{1}{2 P^+}\widetilde \Phi^{[\gamma^+]}_{\unsub} & = 	
		\tAmp_{2B} 
		+ i \mN \epsilon_{ij} \vect{\bvec}_i \vect{S}_j\, \tAmp_{12B} \ ,
\end{align}
where the combined amplitudes are defined as:
\begin{align}
	\tAmp_{2B}  & \ \equiv\  \tAmp_2 + R(\zetahat^2) \tBmp_1 \nonumber \displaybreak[0] \\
	\tAmp_{12B}  & \ \equiv\  \tAmp_{12} - R(\zetahat^2) \tBmp_8 \ . \displaybreak[0] 
	\label{eq-ABamps_sivers}
\end{align}
The final physics deliverables—the unpolarized distribution $f_1$ and the Sivers distribution $f_{1T}^\perp$—are achieved by applying a standardized Fourier integration operator $\fourint$:
\begin{align}
	\fourint \tAmp_i \equiv & \int \frac{d^2 \vprp{\bvec}}{(2\pi)^2}\,e^{-i\vprp{\bvec}\tcdot\vprp{\kei}}\, \frac{1}{\widetilde{\mathcal{S}}(\bvec^2;\ldots)} \int \frac{d(\bvec \tcdot P)}{(2\pi)}\,e^{ix(\bvec \tcdot P)} 
	\tAmp_i(-\vprp{\bvec}^2,\bvec \tcdot P,(\bvec \tcdot P) R(\zetahat^2)/\mN^2,-1/(\mN\zetahat)^2,\eta v \tcdot P) \ .
	\label{eq-fourint}
\end{align}
Applying this to the combined amplitudes reconstructs the target leading-twist distributions:
\begin{align}
	 f_1(x,\vprp{\kei}^2;\zetahat,\ldots,\eta v \tcdot P) & = 2 \fourint\ \widetilde{A}_{2B} \,,\nonumber \displaybreak[0] \\
	 f_{1T}^\perp(x,\vprp{\kei}^2;\zetahat,\ldots,\eta v \tcdot P) & =  4 m_N^2 \partial_{\vprp{\kei}^2} \fourint\ \widetilde{A}_{12B} \ .
	 \label{eq-tmdsfromampsodd}
\end{align} 
Through this framework, the symmetry constraint $f_1(\eta v \tcdot P) = f_1(-\eta v \tcdot P)$ guarantees process independence for the unpolarized distribution, whereas the Sivers function $f_{1T}^\perp$ strictly depends on the direction of the staple link, structurally enforcing the critical sign reversal between SIDIS and DY processes, and laying the groundwork for extracting its $x$-dependence.